\documentclass[12pt, a4paper]{article}

\usepackage[utf8]{inputenc}
\usepackage{amsmath, amssymb, amsthm}
\usepackage{physics}
\usepackage{geometry}
\usepackage{graphicx}
\usepackage{siunitx}
\usepackage{hyperref}

\geometry{margin=1in}

\title{\textbf{Electrodynamic Modeling of a Planar PCB-Embedded Linear Induction Motor (EMALS Analog)}}
\author{Phase 1 Theoretical Framework}
\date{\today}

\begin{document}

\maketitle

\begin{abstract}
This paper establishes the theoretical electrodynamic framework for a solid-state Linear Induction Motor (LIM), miniaturized into the copper layers of an FR4 printed circuit board. Using Maxwell's equations and the magnetic diffusion limit, we derive the time-averaged Lorentz force exerted on a non-magnetic conductive rotor (7075-T6 Aluminum). This analytical model serves as the foundation for the finite element simulation and subsequent custom PCB manufacturing.
\end{abstract}

\section{Introduction}
Modern aerospace launch systems, such as the Electromagnetic Aircraft Launch System (EMALS), rely on linear induction rather than mechanical or steam-driven actuation. To investigate these principles within a SWaP-constrained (Size, Weight, and Power) environment, this project miniaturizes a LIM using a multi-layer Printed Circuit Board (PCB) as the stator. The physical experiment is entirely solid-state, eliminating mechanical bearing friction and allowing pure kinematic analysis of the induced Lorentz forces.

\section{The Traveling Magnetic Field Formulation}
The stator consists of three-phase planar meander coils routed within the inner copper layers of the PCB. When driven by a Sinusoidal Pulse Width Modulation (SPWM) inverter, the $120^\circ$ temporal phase shift and spatial offset $\lambda$ produce a traveling magnetic wave.

We define the propagation axis as $x$ and the orthogonal normal to the PCB plane as $z$. The primary magnetic flux density is defined as:
\begin{equation}
    \mathbf{B}_s(x, t) = B_0 \cos(kx - \omega t)\mathbf{\hat{z}}
\end{equation}
Where $B_0$ is the peak flux density, $k = \frac{\pi}{\tau_p}$ is the wave number (with $\tau_p$ representing the pole pitch), and $\omega = 2\pi f$ is the angular frequency of the applied current. The synchronous phase velocity of this wave is:
\begin{equation}
    v_s = \frac{\omega}{k} = 2f\tau_p
\end{equation}

\section{Rotor Electrodynamics and Slip}
Consider an aluminum rotor of velocity $v = v_x \mathbf{\hat{x}}$. The relative velocity between the traveling wave and the rotor is defined by the dimensionless slip parameter $s$:
\begin{equation}
    s = \frac{v_s - v}{v_s}
\end{equation}
In the rest frame of the rotor, the perceived angular frequency of the changing magnetic field is $\omega_r = s\omega$.

According to Faraday's Law of Induction, $\nabla \times \mathbf{E} = -\frac{\partial \mathbf{B}}{\partial t}$, an electromotive force is induced within the rotor. Assuming a moving conductor in a magnetic field, the generalized Ohm's Law yields the induced eddy current density $\mathbf{J}$:
\begin{equation}
    \mathbf{J} = \sigma(\mathbf{E} + \mathbf{v} \times \mathbf{B})
\end{equation}
Evaluating the cross product for a purely $z$-directed magnetic field and purely $x$-directed velocity, the induced current density is constrained to the transverse $y$-direction:
\begin{equation}
    J_y = \sigma s v_s B_0 \cos(kx - \omega t)
\end{equation}
\textit{Assumption:} We assume the rotor thickness $d$ is significantly less than the magnetic skin depth $\delta = \sqrt{\frac{2}{\mu_0 \sigma \omega_r}}$, allowing us to treat the induced current distribution as uniform across the $z$-axis.

\section{Time-Averaged Lorentz Force Calculation}
The propulsive thrust generated by the LIM is the result of the Lorentz force acting on the induced eddy currents. The volumetric force density $\mathbf{f}$ is:
\begin{equation}
    \mathbf{f} = \mathbf{J} \times \mathbf{B}
\end{equation}
Substituting our derived terms for $\mathbf{J}$ and $\mathbf{B}$:
\begin{align}
    f_x &= J_y B_z \\
    f_x &= \left[ \sigma s v_s B_0 \cos(kx - \omega t) \right] \left[ B_0 \cos(kx - \omega t) \right] \\
    f_x &= \sigma s v_s B_0^2 \cos^2(kx - \omega t)
\end{align}

Because the system operates under high-frequency AC, the instantaneous force oscillates. The physically relevant metric for continuous acceleration is the time-averaged force density $\langle f_x \rangle$ over one spatial period. Integrating $\cos^2(\theta)$ over a full cycle yields a factor of $\frac{1}{2}$:
\begin{equation}
    \langle f_x \rangle = \frac{1}{2} \sigma s v_s B_0^2
\end{equation}
For a rotor of total volume $V$, the net time-averaged propulsive thrust $F_x$ is:
\begin{equation}
    F_x = \frac{1}{2} \sigma V s v_s B_0^2
\end{equation}

\section{Engineering Constraints for PCB Optimization}
Equation (9) mathematically dictates the constraints for the upcoming hardware design phase:
\begin{itemize}
    \item \textbf{Optimal Slip:} Thrust linearly correlates with slip $s$. Force is maximized at $v=0$ ($s=1$) and decays to zero as $v \to v_s$ ($s \to 0$).
    \item \textbf{Flux Dependency:} Thrust scales quadratically with peak magnetic field $B_0$. This requires the PCB trace width to be minimized to maximize spatial turns density, and the MOSFET inverter to be designed with high-current low-$R_{DS(on)}$ silicon to maximize $I_{rms}$ without thermal breakdown of the FR4 substrate.
\end{itemize}

\end{document}